% ═══════════════════════════════════════════════════════════════
% ARBEITSBLATT: Me gusta (Spanisch Oberstufe Spätbeginner)
% Begleitend zum digitalen Lernpfad
% ═══════════════════════════════════════════════════════════════

\documentclass[10pt, a4paper]{article}

% ═══════════════════════════════════════════════════════════════
% PAKETE
% ═══════════════════════════════════════════════════════════════

\usepackage[utf8]{inputenc}
\usepackage[T1]{fontenc}
\usepackage[ngerman, spanish]{babel}
\usepackage{helvet}
\renewcommand{\familydefault}{\sfdefault}

\usepackage[margin=1.5cm, top=1.8cm, bottom=1.5cm]{geometry}
\usepackage{enumitem}
\usepackage{tabularx}
\usepackage{booktabs}
\usepackage{array}
\usepackage{multicol}

\usepackage{xcolor}
\usepackage{tcolorbox}
\usepackage{fancyhdr}
\usepackage{lastpage}
\usepackage[normalem]{ulem}  % für \sout (Durchstreichen)

% Kompaktere Listen
\setlist{nosep, leftmargin=1.5em}

% ═══════════════════════════════════════════════════════════════
% FARBEN
% ═══════════════════════════════════════════════════════════════

\definecolor{spanisch}{HTML}{c41e3a}
\definecolor{grau}{HTML}{888888}
\definecolor{hellgrau}{HTML}{f0f0f0}
\definecolor{success}{HTML}{2a7a4b}

\newcommand{\fachfarbe}{spanisch}

% ═══════════════════════════════════════════════════════════════
% VARIABLEN
% ═══════════════════════════════════════════════════════════════

\newcommand{\thema}{Me gusta}
\newcommand{\fach}{Spanisch}
\newcommand{\klasse}{Oberstufe}

% ═══════════════════════════════════════════════════════════════
% KOPF- UND FUSSZEILE
% ═══════════════════════════════════════════════════════════════

\pagestyle{fancy}
\fancyhf{}
\renewcommand{\headrulewidth}{0.4pt}
\renewcommand{\footrulewidth}{0pt}

\lhead{\small\fach{} \klasse{} -- \thema}
\rhead{\small Nombre: \rule{4cm}{0.4pt}}

\cfoot{\small\thepage/\pageref{LastPage}}

% ═══════════════════════════════════════════════════════════════
% BEFEHLE
% ═══════════════════════════════════════════════════════════════

% Lücke
\newcommand{\luecke}[1]{\rule{#1}{0.4pt}}

% Spanisch-Highlight
\newcommand{\es}[1]{\textcolor{\fachfarbe}{\textbf{#1}}}

% Kompakte Grammatikbox
\newtcolorbox{grammatik}{
    colback=hellgrau,
    colframe=\fachfarbe,
    boxrule=0.8pt,
    arc=0pt,
    left=6pt, right=6pt, top=6pt, bottom=6pt,
    fonttitle=\bfseries\small,
    title=Grammatik
}

% Kompakter Merksatz
\newtcolorbox{merksatz}{
    colback=hellgrau,
    colframe=success,
    boxrule=0.8pt,
    arc=0pt,
    left=6pt, right=6pt, top=4pt, bottom=4pt,
    fonttitle=\bfseries\small,
    title=Recuerda
}

% Fehler-Box
\definecolor{fehler}{HTML}{c62828}
\definecolor{fehlerbg}{HTML}{ffebee}
\newtcolorbox{fehlerbox}{
    colback=fehlerbg,
    colframe=fehler,
    boxrule=0.8pt,
    arc=0pt,
    left=6pt, right=6pt, top=4pt, bottom=4pt,
    fonttitle=\bfseries\small,
    title=Typische Fehler
}

% Abschnittstitel
\newcommand{\abschnitt}[1]{%
    \vspace{8pt}
    \noindent\textbf{\textcolor{\fachfarbe}{#1}}
    \vspace{4pt}
}

% ═══════════════════════════════════════════════════════════════
% DOKUMENT
% ═══════════════════════════════════════════════════════════════

\begin{document}

% ═══════════════════════════════════════════════════════════════
% TITEL
% ═══════════════════════════════════════════════════════════════

\begin{center}
{\large\bfseries Arbeitsblatt: \thema{} -- Hobbys und Vorlieben}\\[2pt]
{\small\textcolor{grau}{Begleitend zum Lernpfad -- in den Hefter übernehmen}}
\end{center}

\vspace{-2pt}
\hrule
\vspace{8pt}

% ═══════════════════════════════════════════════════════════════
% ZWEISPALTIG
% ═══════════════════════════════════════════════════════════════

\begin{multicols}{2}
\small

% ─────────────────────────────────────────────────────────────
% KASTEN 1: Struktur
% ─────────────────────────────────────────────────────────────

\abschnitt{Kasten 1: Gustar = Gefallen}

\begin{grammatik}
\textbf{Wichtig:} Das Verb \es{gustar} funktioniert wie\\
\textbf{„gefallen"}, NICHT wie „mögen"!

\vspace{6pt}

\begin{tabular}{@{}l@{ }l@{}}
\textbf{Deutsch:} & \underline{Mir} gefällt \underline{die Musik}. \\[2pt]
\textbf{Español:} & \es{Me} gusta \es{la música}. \\
\end{tabular}

\vspace{6pt}

\textbf{Struktur:}\\
Wem? $\rightarrow$ Indirektes Objektpronomen\\
Was? $\rightarrow$ Subjekt (bestimmt die Verbform!)
\end{grammatik}

\vspace{4pt}

\begin{fehlerbox}
\begin{tabular}{@{}l@{ }l@{}}
\textcolor{fehler}{\sout{*Yo gusto música.}} & $\rightarrow$ Falsche Struktur! \\
\textcolor{fehler}{\sout{*Me gusta los libros.}} & $\rightarrow$ Plural: \es{gustan}! \\
\textcolor{fehler}{\sout{*Me gustan bailar.}} & $\rightarrow$ Infinitiv: \es{gusta}! \\
\end{tabular}
\end{fehlerbox}

\vspace{6pt}

% ─────────────────────────────────────────────────────────────
% KASTEN 2: Pronomen
% ─────────────────────────────────────────────────────────────

\abschnitt{Kasten 2: Los pronombres}

\begin{tabular}{|l|l|l|}
\hline
\textbf{Person} & \textbf{Pronomen} & \textbf{Deutsch} \\
\hline
(A mí) & \es{me} & mir \\
\hline
(A ti) & \es{te} & dir \\
\hline
(A él/ella) & \es{le} & ihm/ihr \\
\hline
(A nosotros) & \es{nos} & uns \\
\hline
(A vosotros) & \es{os} & euch \\
\hline
(A ellos) & \es{les} & ihnen \\
\hline
\end{tabular}

\vspace{6pt}

\begin{merksatz}
\es{gusta} + Singular / Infinitiv\\
\es{gustan} + Plural
\end{merksatz}

\vspace{6pt}

{\small\textbf{Ejemplos:}}

{\small
\begin{tabular}{@{}l@{ }l@{}}
\es{Me gusta bailar.} & Mir gefällt tanzen. \\
\es{Te gustan los libros.} & Dir gefallen die Bücher. \\
\es{A María le gusta el café.} & María gefällt Kaffee. \\
\es{No nos gusta cocinar.} & Uns gefällt Kochen nicht. \\
\end{tabular}}

\vspace{4pt}

{\footnotesize\textbf{Wozu „A mí", „A ti"...?}

\textit{Optional} -- aber nützlich für:
\begin{itemize}[nosep, leftmargin=1em]
\item \textbf{Betonung:} \es{A mí} me gusta, \es{a ti} no.
\item \textbf{Klarstellung bei le/les:} \es{A María} le gusta...
\end{itemize}}

\vspace{6pt}

% ─────────────────────────────────────────────────────────────
% VOKABELN
% ─────────────────────────────────────────────────────────────

\abschnitt{Vocabulario: Hobbys y tiempo libre}

\begin{tabular}{|l|l|}
\hline
\textbf{Español} & \textbf{Deutsch} \\
\hline
la música & die Musik \\
\hline
el deporte & der Sport \\
\hline
leer & lesen \\
\hline
bailar & tanzen \\
\hline
cocinar & kochen \\
\hline
viajar & reisen \\
\hline
jugar al fútbol & Fußball spielen \\
\hline
ver películas & Filme schauen \\
\hline
escuchar música & Musik hören \\
\hline
salir con amigos & mit Freunden ausgehen \\
\hline
\end{tabular}

% ═══════════════════════════════════════════════════════════
% SPALTE 2
% ═══════════════════════════════════════════════════════════

\columnbreak

% ─────────────────────────────────────────────────────────────
% ÜBUNG 1
% ─────────────────────────────────────────────────────────────

\abschnitt{Übung 1: Pronomen einsetzen}

Setze das richtige Pronomen ein (me, te, le, nos, os, les):

\vspace{4pt}

\begin{enumerate}[leftmargin=2em]
\item \luecke{1.2cm} gusta el chocolate. (mir)
\item \luecke{1.2cm} gustan los videojuegos. (dir)
\item A María \luecke{1.2cm} gusta leer. (ihr)
\item \luecke{1.2cm} gusta viajar. (uns)
\item A Pedro y Ana \luecke{1.2cm} gustan los gatos. (ihnen)
\end{enumerate}

\vspace{8pt}

% ─────────────────────────────────────────────────────────────
% ÜBUNG 2
% ─────────────────────────────────────────────────────────────

\abschnitt{Übung 2: gusta oder gustan?}

Ergänze die richtige Verbform:

\vspace{4pt}

\begin{enumerate}[leftmargin=2em]
\item Me \luecke{1.5cm} la música.
\item Te \luecke{1.5cm} los deportes.
\item Le \luecke{1.5cm} bailar.
\item Nos \luecke{1.5cm} las películas.
\item Les \luecke{1.5cm} el fútbol.
\end{enumerate}

\vspace{8pt}

% ─────────────────────────────────────────────────────────────
% ÜBUNG 3
% ─────────────────────────────────────────────────────────────

\abschnitt{Übung 3: Übersetze ins Spanische}

\vspace{4pt}

\begin{enumerate}[leftmargin=2em]
\item Mir gefällt Musik. \luecke{4cm}
\item Dir gefallen die Filme. \luecke{4cm}
\item Uns gefällt das Reisen. \luecke{4cm}
\end{enumerate}

\vspace{6pt}

% ─────────────────────────────────────────────────────────────
% ÜBUNG 4: BEGRÜNDEN (AFB III)
% ─────────────────────────────────────────────────────────────

\abschnitt{Übung 4: Fehler erklären}

Warum sind diese Sätze falsch? Erkläre kurz:

\vspace{4pt}

\begin{enumerate}[leftmargin=2em]
\item \textcolor{fehler}{\sout{*Me gustan bailar.}}

\luecke{5cm}

\vspace{2pt}

\item \textcolor{fehler}{\sout{*Yo gusto el chocolate.}}

\luecke{5cm}

\end{enumerate}

\vspace{6pt}

% ─────────────────────────────────────────────────────────────
% FREIES SCHREIBEN
% ─────────────────────────────────────────────────────────────

\abschnitt{Freies Schreiben}

Schreibe 3 Sätze: Was magst du? Was magst du nicht?

\vspace{4pt}

1. \luecke{5.5cm}

\vspace{6pt}

2. \luecke{5.5cm}

\vspace{6pt}

3. \luecke{5.5cm}

\vspace{8pt}

% ─────────────────────────────────────────────────────────────
% LÖSUNGEN
% ─────────────────────────────────────────────────────────────

{\scriptsize\textcolor{grau}{
\textbf{Lösungen:}
Ü1: me, te, le, nos, les |
Ü2: gusta, gustan, gusta, gustan, gusta |
Ü3: Me gusta la música. / Te gustan las películas. / Nos gusta viajar. |
Ü4: 1) Infinitiv braucht gusta (Sg.) 2) gustar = gefallen, nicht mögen
}}

\end{multicols}

\end{document}
